\documentclass[12pt]{article}
\usepackage[utf8]{inputenc}
\usepackage{amsmath, amssymb, amsthm}
\usepackage{geometry}
\usepackage{fancyhdr}

\geometry{letterpaper, margin=1in}
\setlength{\headheight}{15pt}

% Header setup
\pagestyle{fancy}
\fancyhf{}
\lhead{Math 477 Spring 2026}
\rhead{Jeffrey Lu}
\cfoot{\thepage}

\title{Homework Assignment 1 \\ \large Math 292, Spring 2026 \\ Rutgers University}
\author{Jeffrey Lu}

\begin{document}

\maketitle

\noindent \textbf{Question 4.8}

Suppose $a$ is an integer. Assuming $5 ~|~ 2a$, there exists an integer k such that
\[
    2a = 5k.
\]

Multiplying both sides by 3 and rearranging we have that
\begin{align*}
    6a & = 15k                           \\
    a  & = 15k -5a                       \\
    a  & = 5(3k-a)                       \\
    a  & = 5w ~|~ \text{where w = 3k-a.}
\end{align*}

And, w must be an integer because it is the linear combination of integers.
Thus, by the definition of divisibility, we have
\[
    5 ~|~ a.
\]

Therefore, we can conclude that if $ 5 ~|~ 2a$, then $5 ~|~ a$. $\square{}$

\vspace{50pt}

\noindent \textbf{Question 4.12}

Suppose $x \in{\mathbb{R}}$ and $0 < x < 4$.
We let a = (x-2), where a is a real number because it is the linear combination of real numbers.
Because the square of a real number must be nonnegative, we have that
\begin{align*}
    a^2 = (x-2)^2 \geq{0}.
\end{align*}
Expanding this we get
\begin{align*}
    x^2-4x+4 \geq{0}.
\end{align*}
Then, we divide through by -1, flipping the inequality and resulting in
\begin{align*}
    -x^2+4x-4 & \leq{0} \\
    -x^2+4x   & \leq{4} \\
    x(4-x)    & \leq{4} \\
\end{align*}
Recall that $0<x<4$, which gives us these two facts

1. $x > 0$

2. $4 - x > 0$

\noindent We let $b = 4-x$ and $c = x*b = x(4-x)$. Since b and x are both positive, c being the product of two positive real numbers must also be positive.

Thus, we can safely divide through by x(4-x), giving us
\begin{align*}
    1 \leq{\frac{4}{x(4-x)}}.
\end{align*}

Therefore we conclude that if $0 < x < 4$, then $\frac{4}{x(4 - x)} \geq 1$. $\square{}$

\vspace{50pt}

\noindent \textbf{Question 4.14}

Suppose $n \in{\mathbb{Z}}$. Then, we have the two possible cases.
$n$ could either be a odd integer or a even integer.
\vspace{4pt}

Case 1: n is even. If n is an even integer, we can rewrite it as $n = 2x$ for some integer x.
Substituting this back into our expression we now have
\begin{align*}
    5n^2 +3n + 7 & = 5(2x^2) + 3(2x) + 7    \\
                 & = 20x^2 + 6x + 7         \\
                 & = 20x^2 + 6x + 6 + 1     \\
                 & = 2(10x^2 + 3x + 3) + 1.
\end{align*}

Let $j = 10x^2 + 3x + 3$. Because j is the linear combination of integers, j is also an integer.


Hence, our expression reduces to
\[
    5n^2 +3n + 7 = 2j+1 \text{, where } j\in{\mathbb{Z}}
\]

Therefore, by definition of odd, we can conclude that $5n^2 +3n + 7$ is odd.
\vspace{8pt}

Case 2: n is odd. If n is an odd integer, we can rewrite it as $n = 2x+1$ for some integer x.
Substituting this back into our expression we now have
\begin{align*}
    5n^2 +3n + 7 & = 5(2x+1)^2 + 3(2x+1) + 7      \\
                 & = 20x^2 + 20x + 5 + 6x + 3 + 7 \\
                 & = 20x^2 + 26x + 12 + 1         \\
                 & = 2(10x^2 + 13x + 6) + 1.
\end{align*}

Let $j = 10x^2 + 13x + 6$. Because j is the linear combination of integers, j is also an integer.

Hence, our expression reduces to
\[
    5n^2 +3n + 7 = 2j+1 \text{, where } j\in{\mathbb{Z}}.
\]
Therefore, by definition of odd, we can conclude that $5n^2 +3n + 7$ is odd.
\vspace{8pt}

In both cases, we have shown that the expression turns out to be odd.
Therefore, we conclude that if $n\in{\mathbb{Z}}$, then $5n^2 +3n + 7$ is odd. $\square$

\vspace{50pt}

\noindent \textbf{Question 5.6}

Using the contrapositive, we have that given a real number x, suppose that $x \leq{-1}$, then we must prove that $x^3-x \leq{0}$.

Adding and subtracting 1 from x, we get
\begin{align*}
    x - 1 & \leq{-2} \\
    x + 1 & \leq{0}.
\end{align*}

Multiplying x by x-1, which is strictly negative, flips the sign, giving us
\[
    x(x-1) \geq{2}.
\]

Then, multiplying by x + 1, which is non-positive, we flip the sign again
\[
    x(x-1)(x+1) \leq{0}.
\]

Combining the terms, we have
\[
    x^3 -x \leq{0}.
\]

Therefore, we conclude that if if $x^3 - x > 0$, then $x > -1$. $\Box$

\vspace{50pt}

\noindent \textbf{Question 18a}

Suppose $x < y$. Rearranging, we get
\[
    y-x > 0
\]

Now, consider $x^2+xy+y^2$, completing the square, we have
\begin{align*}
    x^2+xy+y^2 = (x+\frac{y}{2})^2 + \frac{3}{4}y^2
\end{align*}

Let $a = (x+\frac{y}{2})$.

We know a must be a real number since any linear combination of reals is still a real number.
Thus, we have
\[
    x^2+xy+y^2 = a^2 + \frac{3}{4}y^2.
\]

Any real number squared must be non-negative and we have that x and y can't both be zero from the original inequality.
Then, we have two cases either $a = 0$ then $y^2>0$, either $a \neq{0}$ then $a^2 = 0$, thus we have
\[
    x^2+xy+y^2 = a^2 + \frac{3}{4}y^2 > 0.
\]

Next, we multiply both sides by $(y-x)$, which we know is strictly positive and won't flip the inequality
\[
    (y-x)(x^2+xy+y^2) > 0
\]

Distributing the terms, we have
\begin{align*}
    y(y^2 + xy + x^2) - x(y^2 + xy + x^2) & > 0 \\
    y^3 + xy^2 + x^2y - xy^2 - x^2y - x^3 & > 0 \\
    y^3 - x^3                             & > 0
\end{align*}

Adding $x^3$ on both sides, we get
\[
    y^3 > x^3
\]

Therefore, we can conclude that if $x< y$, then $x^3 < y^3$. $\Box{}$

\vspace{50pt}

\noindent \textbf{Question 29a}

Let x be a real number. Let a = x -1.
Then, a must be a real number because it is the linear combination of real numbers.
Since the square of a real number is non-negative we have
\begin{align*}
    a^2 = (x-1)^2 \geq{0}.
\end{align*}

Expanding, we get
\[
    x^2-2x+1 \geq{0}.
\]

Then, we add $6x+3$ to both sides
\[
    x^2+4x+4 \geq{6x+3}.
\]

We can now factor the left expression
\[
    (x+2)^2\geq{6x+3}
\]

Therefore, we can conclude that for any real number x, we have that $(x+2)^2\geq{6x+3} \Box{}$

\end{document}